\section*{Chapter \ref{ch:b3:human-capital} --- Human Capital and the Returns to Schooling}

\begin{solution}{exo:b3:human-capital:1}
Lemieux and Card's Canadian estimate: 16.4~\% by IV against 7.0~\% by OLS, a gap of 9.4 points,
just above Harmon and Walker's 15.3 against 6.1. Duflo's Indonesian estimate is the only
one below: 6.4 against 7.8. % figdata: estimates
\end{solution}

\begin{solution}{exo:b3:human-capital:2}
$1{,}500/20{,}000 = 7.5~\%$: the rate at which a perpetuity of 1{,}500 is worth 20{,}000.
\end{solution}

\begin{solution}{exo:b3:human-capital:3}
$e^{0.11} - 1 = 11.6~\%$.
\end{solution}

\begin{solution}{exo:b3:human-capital:4}
A return within pairs of identical twins, corrected for reporting error, clearly
below the OLS return across people.
\end{solution}

\begin{solution}{exo:b3:human-capital:5}
The tested solution is \code{ex05_mincer.py}, in the chapter's \code{exercises/}
folder. The coefficient falls from 0.0632 (s.e.\ 0.0032) to 0.0552 (0.0033): a premium of 6.5 and
then 5.7~\% a year. % output: ex05_mincer
The added indicators are correlated with both schooling and earnings, and absorb part
of what the first regression credited to schooling.
\end{solution}

\begin{solution}{exo:b3:human-capital:6}
Twins' schooling is highly correlated, so the true variance of the difference is small,
while the variance of the reporting error is doubled by differencing: the share of
noise in the measured difference, and so the bias toward zero, is larger.
\end{solution}

\begin{solution}{exo:b3:human-capital:7}
$0.024/0.3 = 0.08$: 8~\% per year.
\end{solution}

\begin{solution}{exo:b3:human-capital:8}
A forklift licence is general: every warehouse values it, so the worker pays,
directly or through a lower wage while she trains. Training in the firm's own
software is specific: the firm can recover its cost and shares it.
\end{solution}

\begin{solution}{exo:b3:human-capital:9}
A coefficient on the diploma indicator of zero: each year adds the same, whether it
completes a diploma or not.
\end{solution}

\begin{solution}{exo:b3:human-capital:10}
The proof is in the text. With tuition $c$ paid at the start of year $S + 1$, equal value requires
$Y(S+1)e^{-r(S+1)}/r = Y(S)e^{-rS}/r + ce^{-rS}$, so $Y(S+1)e^{-r} = Y(S) + rc$ and $\ln Y(S+1) - \ln Y(S) = r + \ln(1 + rc/Y(S)) \approx r(1 + c/Y(S))$.
\end{solution}

\begin{solution}{exo:b3:human-capital:11}
The tested solution is \code{ex11_profiles.py}, in the chapter's \code{exercises/}
folder. The fitted profiles peak at 43.6 years for men with fewer than 12 years of
school, 41.0 with 12, 40.6 with 13--15 and 39.7 with 16 or more. % output: ex11_profiles
Top-coding caps the earnings of older, more educated men, flattening their profiles
where they would still rise and pulling the fitted peak earlier.
\end{solution}

\begin{solution}{exo:b3:human-capital:12}
IV estimates the return of the \omterm{def:b3:instrumental-variables:compliance}{compliers} --- those whose schooling the instrument
moves. If they are people kept from school by its cost or distance, their marginal
return can be well above the average, and above the OLS estimate even when ability
biases OLS upward.
\end{solution}

\begin{solution}{pb:b3:human-capital:school}
% figdata: brenland_stats
\textbf{1.} 1{,}439 heads; mean schooling 12.58 years.

\textbf{2.} 0.083.

\textbf{3.} 0.081 (s.e.\ 0.004).

\textbf{4.} $e^{0.081} - 1 = 8.4~\%$.

\textbf{5.} Upward, if the able both study longer and earn more.

\textbf{6.} The survey gives the head's gross income, not an hourly wage; the slope then
includes any effect of schooling on hours and on income from savings.

\textbf{7.} 20.4~\%.

\textbf{8.} 8.99 years; the reform adds $11 - 8.99 = 2.01$ years to each.

\textbf{9.} $0.204 \times 2.01 = 0.411$ years.

\textbf{10.} $0.029/0.411 = 7.06~\%$.

\textbf{11.} It is below the least-squares 8.1~\%, as upward \omterm{def:b3:human-capital:ability-bias}{ability bias} would make it --- or
the \omterm{def:b3:instrumental-variables:compliance}{compliers}' return is lower than the average.

\textbf{12.} The reform changes earnings only through schooling. It fails if the
reformed cohorts entered a different labour market, or if the reform came with
other changes to schools.

\textbf{13.} Those who would have left with 10 years or fewer.

\textbf{14.} They are the least schooled: their return may be lower (less interest or
preparation) or higher (kept out by the cost of staying).

\textbf{15.} That they were held back by constraints rather than by a low return.

\textbf{16.} It moves only the schooling of those at the 10--11 margin.

\textbf{17.} Distance to a university as an instrument, or a regression discontinuity
at an admission threshold.

\textbf{18.} After tax, earnings are $(1 - t)Y$ in both streams and the cost is $(1 - t)Y$ given up: the
log gain and the ratio of gain to cost are unchanged, so the return is 7.06~\%.

\textbf{19.} $0.75 \times 7.06 = 5.30~\%$.

\textbf{20.} $5.30/1.3 = 4.07~\%$.

\textbf{21.} Taxes move income between the pupil and the state; output is counted
before tax.

\textbf{22.} Yes: the year costs the same and adds less output.

\textbf{23.} Gains to others: less crime, more informed voters, colleagues who learn
from the better educated.

\textbf{24.} 4.07~\% exceeds 3~\%: at that discount rate the year is worth its cost to society.

\textbf{25.} The private return is 7.1~\% and the social return 4.1~\%.
\end{solution}
